% Image credits for the photographs and AI illustrations of Book 4
% (University Biology, Year 2). Machine-readable license records live in
% images/book4/CREDITS.md; keep the two files in sync.
\chapter*{Crédits des images}

% \sloppy must be in force when the paragraph is BROKEN, hence the \par before
% \endgroup -- without it the group closes first and the tolerance is restored.
\begingroup\sloppy

Les dessins de ce livre sont originaux. Les photographies sont
utilisées sous leurs licences libres respectives, avec nos
remerciements à leurs auteurs~:

\begin{itemize}
  \item Chapitre~\ref{ch:b2:microbial-metabolism}~: photographie de
        Sergueï Winogradsky, années 1890 --- domaine public (Wikimedia
        Commons).
  \item Chapitre~\ref{ch:b2:genome-diversification}~: photographie de
        Barbara McClintock dans son laboratoire, Smithsonian Institution
        / Science Service, 1947 --- domaine public (Wikimedia Commons).
  \item Chapitre~\ref{ch:b2:blood-circulation}~: portrait de William
        Harvey, et planche des veines de l'avant-bras tirée de son
        \emph{Exercitatio anatomica de motu cordis} (1628), Wellcome
        Collection --- CC~BY~4.0 (Wikimedia Commons)~; micrographie
        électronique à balayage d'une hématie, d'une plaquette et d'un
        lymphocyte, Electron Microscopy Facility, National Cancer
        Institute --- domaine public (Wikimedia Commons).
  \item Chapitre~\ref{ch:b2:speciation}~: planche des pinsons de Darwin
        par John Gould, tirée du \emph{Journal of Researches} (1845) ---
        domaine public (Wikimedia Commons).
  \item Chapitre~\ref{ch:b2:biogeochemical-cycles}~: l'observatoire du
        Mauna Loa (Forrest M. Mims III, 2006) et un prélèvement d'air en
        flacon effectué sur place (Commander John Bortniak, 1982), NOAA
        Photo Library --- domaine public (Wikimedia Commons).
\end{itemize}

Les liens vers les sources complètes et les adresses des licences de
chaque photographie sont donnés dans
\texttt{images/\allowbreak book4/\allowbreak CREDITS.md}, dans le dépôt du livre.

\bigskip
À côté des photographies et des dessins originaux, environ
soixante-dix figures réparties dans les chapitres du volume
sont des illustrations produites par intelligence artificielle,
réalisées avec la génération d'images d'OpenAI à partir d'invites
écrites par les éditeurs du livre, et relues pour en vérifier
l'exactitude biologique avant leur inclusion. L'invite complète de chaque image générée est
consignée dans \texttt{images/\allowbreak book4/\allowbreak ai/\allowbreak PROMPTS.md},
dans le dépôt.
\par\endgroup
\clearpage
